%&preformat-disser
\RequirePackage[l2tabu,orthodox]{nag} % Раскомментировав, можно в логе получать рекомендации относительно правильного использования пакетов и предупреждения об устаревших и нерекомендуемых пакетах
% Формат А4, 14pt (ГОСТ Р 7.0.11-2011, 5.3.6)
\documentclass[a4paper,14pt,oneside,openany]{memoir}

\input{common/setup}            % общие настройки шаблона
\input{common/packages}         % Пакеты общие для диссертации и автореферата
\synopsisfalse                      % Этот документ --- не автореферат
\input{Dissertation/dispackages}    % Пакеты для диссертации
\input{Dissertation/userpackages}   % Пакеты для специфических пользовательских задач

\input{Dissertation/setup}      % Упрощённые настройки шаблона

\input{common/newnames}         % Новые переменные, для всего проекта

%%% Основные сведения %%%
\newcommand{\thesisAuthorLastName}{\fixme{Ступников}}
\newcommand{\thesisAuthorOtherNames}{\fixme{Вадим Владимирович}}
\newcommand{\thesisAuthorInitials}{\fixme{И.\,О.}}
\newcommand{\thesisAuthor}             % Диссертация, ФИО автора
{%
    \texorpdfstring{% \texorpdfstring takes two arguments and uses the first for (La)TeX and the second for pdf
        \thesisAuthorLastName~\thesisAuthorOtherNames% так будет отображаться на титульном листе или в тексте, где будет использоваться переменная
    }{%
        \thesisAuthorLastName, \thesisAuthorOtherNames% эта запись для свойств pdf-файла. В таком виде, если pdf будет обработан программами для сбора библиографических сведений, будет правильно представлена фамилия.
    }
}
\newcommand{\thesisAuthorShort}        % Диссертация, ФИО автора инициалами
{\thesisAuthorInitials~\thesisAuthorLastName}
%\newcommand{\thesisUdk}                % Диссертация, УДК
%{\fixme{xxx.xxx}}
\newcommand{\thesisTitle}              % Диссертация, название
{\fixme{Длинное название диссертационной работы, состоящее из~достаточно большого
количества слов, совсем длинное длинное длинное длинное название, из~которого
простому обывателю знакомы, в~лучшем случае, лишь отдельные слова}}
\newcommand{\thesisSpecialtyNumber}    % Диссертация, специальность, номер
{\fixme{XX.XX.XX}}
\newcommand{\thesisSpecialtyTitle}     % Диссертация, специальность, название (название взято с сайта ВАК для примера)
{\fixme{Технология обработки, хранения и~переработки злаковых, бобовых культур,
крупяных продуктов, плодоовощной продукции и~виноградарства}}
%% \newcommand{\thesisSpecialtyTwoNumber} % Диссертация, вторая специальность, номер
%% {\fixme{XX.XX.XX}}
%% \newcommand{\thesisSpecialtyTwoTitle}  % Диссертация, вторая специальность, название
%% {\fixme{Теория и~методика физического воспитания, спортивной тренировки,
%% оздоровительной и~адаптивной физической культуры}}
\newcommand{\thesisDegree}             % Диссертация, ученая степень
{\fixme{кандидата физико-математических наук}}
\newcommand{\thesisDegreeShort}        % Диссертация, ученая степень, краткая запись
{\fixme{канд. физ.-мат. наук}}
\newcommand{\thesisCity}               % Диссертация, город написания диссертации
{\fixme{Город}}
\newcommand{\thesisYear}               % Диссертация, год написания диссертации
{\the\year}
\newcommand{\thesisOrganization}       % Диссертация, организация
{\fixme{Федеральное государственное автономное образовательное учреждение высшего
образования <<Длинное название образовательного учреждения <<АББРЕВИАТУРА>>}}
\newcommand{\thesisOrganizationShort}  % Диссертация, краткое название организации для доклада
{\fixme{НазУчДисРаб}}

\newcommand{\thesisInOrganization}     % Диссертация, организация в предложном падеже: Работа выполнена в ...
{\fixme{учреждении с~длинным длинным длинным длинным названием, в~котором
выполнялась данная диссертационная работа}}

%% \newcommand{\supervisorDead}{}           % Рисовать рамку вокруг фамилии
\newcommand{\supervisorFio}              % Научный руководитель, ФИО
{\fixme{Фамилия Имя Отчество}}
\newcommand{\supervisorRegalia}          % Научный руководитель, регалии
{\fixme{уч. степень, уч. звание}}
\newcommand{\supervisorFioShort}         % Научный руководитель, ФИО
{\fixme{И.\,О.~Фамилия}}
\newcommand{\supervisorRegaliaShort}     % Научный руководитель, регалии
{\fixme{уч.~ст.,~уч.~зв.}}

%% \newcommand{\supervisorTwoDead}{}        % Рисовать рамку вокруг фамилии
%% \newcommand{\supervisorTwoFio}           % Второй научный руководитель, ФИО
%% {\fixme{Фамилия Имя Отчество}}
%% \newcommand{\supervisorTwoRegalia}       % Второй научный руководитель, регалии
%% {\fixme{уч. степень, уч. звание}}
%% \newcommand{\supervisorTwoFioShort}      % Второй научный руководитель, ФИО
%% {\fixme{И.\,О.~Фамилия}}
%% \newcommand{\supervisorTwoRegaliaShort}  % Второй научный руководитель, регалии
%% {\fixme{уч.~ст.,~уч.~зв.}}

\newcommand{\opponentOneFio}           % Оппонент 1, ФИО
{\fixme{Фамилия Имя Отчество}}
\newcommand{\opponentOneRegalia}       % Оппонент 1, регалии
{\fixme{доктор физико-математических наук, профессор}}
\newcommand{\opponentOneJobPlace}      % Оппонент 1, место работы
{\fixme{Не очень длинное название для места работы}}
\newcommand{\opponentOneJobPost}       % Оппонент 1, должность
{\fixme{старший научный сотрудник}}

\newcommand{\opponentTwoFio}           % Оппонент 2, ФИО
{\fixme{Фамилия Имя Отчество}}
\newcommand{\opponentTwoRegalia}       % Оппонент 2, регалии
{\fixme{кандидат физико-математических наук}}
\newcommand{\opponentTwoJobPlace}      % Оппонент 2, место работы
{\fixme{Основное место работы c длинным длинным длинным длинным названием}}
\newcommand{\opponentTwoJobPost}       % Оппонент 2, должность
{\fixme{старший научный сотрудник}}

%% \newcommand{\opponentThreeFio}         % Оппонент 3, ФИО
%% {\fixme{Фамилия Имя Отчество}}
%% \newcommand{\opponentThreeRegalia}     % Оппонент 3, регалии
%% {\fixme{кандидат физико-математических наук}}
%% \newcommand{\opponentThreeJobPlace}    % Оппонент 3, место работы
%% {\fixme{Основное место работы c длинным длинным длинным длинным названием}}
%% \newcommand{\opponentThreeJobPost}     % Оппонент 3, должность
%% {\fixme{старший научный сотрудник}}

\newcommand{\leadingOrganizationTitle} % Ведущая организация, дополнительные строки. Удалить, чтобы не отображать в автореферате
{\fixme{Федеральное государственное бюджетное образовательное учреждение высшего
профессионального образования с~длинным длинным длинным длинным названием}}

\newcommand{\defenseDate}              % Защита, дата
{\fixme{DD mmmmmmmm YYYY~г.~в~XX часов}}
\newcommand{\defenseCouncilNumber}     % Защита, номер диссертационного совета
{\fixme{Д\,123.456.78}}
\newcommand{\defenseCouncilTitle}      % Защита, учреждение диссертационного совета
{\fixme{Название учреждения}}
\newcommand{\defenseCouncilAddress}    % Защита, адрес учреждение диссертационного совета
{\fixme{Адрес}}
\newcommand{\defenseCouncilPhone}      % Телефон для справок
{\fixme{+7~(0000)~00-00-00}}

\newcommand{\defenseSecretaryFio}      % Секретарь диссертационного совета, ФИО
{\fixme{Фамилия Имя Отчество}}
\newcommand{\defenseSecretaryRegalia}  % Секретарь диссертационного совета, регалии
{\fixme{д-р~физ.-мат. наук}}            % Для сокращений есть ГОСТы, например: ГОСТ Р 7.0.12-2011 + http://base.garant.ru/179724/#block_30000

\newcommand{\synopsisLibrary}          % Автореферат, название библиотеки
{\fixme{Название библиотеки}}
\newcommand{\synopsisDate}             % Автореферат, дата рассылки
{\fixme{DD mmmmmmmm}\the\year~года}

% To avoid conflict with beamer class use \providecommand
\providecommand{\keywords}%            % Ключевые слова для метаданных PDF диссертации и автореферата
{}
             % Основные сведения
\input{common/fonts}            % Определение шрифтов (частичное)
\input{common/styles}           % Стили общие для диссертации и автореферата
\input{Dissertation/disstyles}  % Стили для диссертации
\input{Dissertation/userstyles} % Стили для специфических пользовательских задач

%%% Библиография. Выбор движка для реализации %%%
% Здесь только проверка установленного ключа. Сама настройка выбора движка
% размещена в common/setup.tex
\ifnumequal{\value{bibliosel}}{0}{%
    \input{biblio/predefined}   % Встроенная реализация с загрузкой файла через движок bibtex8
}{
    \input{biblio/biblatex}     % Реализация пакетом biblatex через движок biber
}

% Вывести информацию о выбранных опциях в лог сборки
\typeout{Selected options:}
\typeout{Draft mode: \arabic{draft}}
\typeout{Font: \arabic{fontfamily}}
\typeout{AltFont: \arabic{usealtfont}}
\typeout{Bibliography backend: \arabic{bibliosel}}
\typeout{Precompile images: \arabic{imgprecompile}}
% Вывести информацию о версиях используемых библиотек в лог сборки
\listfiles

%%% Управление компиляцией отдельных частей диссертации %%%
% Необходимо сначала иметь полностью скомпилированный документ, чтобы все
% промежуточные файлы были в наличии
% Затем, для вывода отдельных частей можно воспользоваться командой \includeonly
% Ниже примеры использования команды:
%
%\includeonly{Dissertation/part2}
%\includeonly{Dissertation/contents,Dissertation/appendix,Dissertation/conclusion}
%
% Если все команды закомментированы, то документ будет выведен в PDF файл полностью

\begin{document}
%%% Переопределение именований типовых разделов
% https://tex.stackexchange.com/a/156050
\gappto\captionsrussian{\input{common/renames}\unskip} % for polyglossia and babel
\input{common/renames}

%%% Структура диссертации (ГОСТ Р 7.0.11-2011, 4)
\thispagestyle{empty}

\noindent%
\begin{tabularx}{\textwidth}{@{}lXr@{}}%
    & & \large{На правах рукописи}\\
    \IfFileExists{images/logo.png}{\includegraphics[height=2.0cm]{logo}}{\rule[0pt]{0pt}{2.0cm}}  & &
    \ifnumequal{\value{showperssign}}{0}{%
        \rule[0pt]{0pt}{1.5cm}
    }{
        \includegraphics[height=1.5cm]{personal-signature.png}
    }\\
\end{tabularx}

\vspace{0pt plus0fill} %число перед fill = кратность относительно некоторого расстояния fill, кусками которого заполнены пустые места
\begin{center}
%\textbf {\large \thesisAuthor} УДАЛИЛ ПОД ЕВСЮКОВА
\textbf {\thesisAuthor}
\end{center}

\vspace{0pt plus2.5fill} %число перед fill = кратность относительно некоторого расстояния fill, кусками которого заполнены пустые места
\begin{center}
%\textbf {\Large %\MakeUppercase	УДАЛИЛ ПОД ЕВСЮКОВА
%\thesisTitle}						УДАЛИЛ ПОД ЕВСЮКОВА
\textbf { \MakeUppercase
	\thesisTitle}

\vspace{0pt plus5.0fill} %число перед fill = кратность относительно некоторого расстояния fill, кусками которого заполнены пустые места
{\textbf{\thesisSpecialtyNumber\ \thesisSpecialtyTitle}}
%{\large Специальность \thesisSpecialtyNumber\ "---\par <<\thesisSpecialtyTitle>>}

\ifdefined\thesisSpecialtyTwoNumber
{\large Специальность \thesisSpecialtyTwoNumber\ "---\par <<\thesisSpecialtyTwoTitle>>}
\fi

\vspace{0pt plus8.0fill} %число перед fill = кратность относительно некоторого расстояния fill, кусками которого заполнены пустые места
Автореферат диссертации на соискание учёной степени\par
\thesisDegree


%\Large{Автореферат}\par
%\large{диссертации на соискание учёной степени\par \thesisDegree}
\end{center}

\vspace{0pt plus4fill} %число перед fill = кратность относительно некоторого расстояния fill, кусками которого заполнены пустые места
%{\centering\thesisCity~--- \thesisYear\par} УДАЛИЛ ПОД ЕВСЮКОВА
{\centering\textbf{\thesisCity~\thesisYear}\par}

\newpage
% оборотная сторона обложки
\thispagestyle{empty}
\noindent Работа выполнена в {\thesisInOrganizationFull}.

\vspace{0.008\paperheight plus1fill}
\noindent%
\begin{tabularx}{\textwidth}{@{}lX@{}}
    \ifdefined\supervisorTwoFio
    Научные руководители:   & \supervisorRegalia\par
                              \ifdefined\supervisorDead
                              \framebox{\textbf{\supervisorFio}}
                              \else
                              \textbf{\supervisorFio}
                              \fi
                              \par
                              \vspace{0.013\paperheight}
                              \supervisorRegalia\par
                              \ifdefined\supervisorTwoDead
                              \framebox{\textbf{\supervisorTwoFio}}
                              \else
                              \textbf{\supervisorTwoFio}
                              \fi
                              \vspace{0.013\paperheight}\\
    \else
    \textbf{Научный руководитель:}   & \textbf{\supervisorFio}\par
                              \ifdefined\supervisorDead
                              \framebox{\textbf{\supervisorFio}}
                              \else
                              \supervisorRegalia
                              \fi
                              \vspace{0.013\paperheight}\\
    \fi
    \textbf{Официальные оппоненты:}  &
    \ifnumequal{\value{showopplead}}{0}{\vspace{13\onelineskip plus1fill}}{%
        \textbf{\opponentOneFio}\par
% Ниже сделаны правки под Евсюкова
%\textbf{\opponentOneFio,}\par
        \opponentOneRegalia\par
%\opponentOneRegalia,\par
%        \opponentOneJobPlace\par
%\opponentOneJobPlace,\par
%        \opponentOneJobPost\par
        \vspace{0.01\paperheight}
        \textbf{\opponentTwoFio}\par
        \opponentTwoRegalia\par
%        \opponentTwoJobPlace\par
%\textbf{\opponentTwoFio,}\par
%\opponentTwoRegalia,\par
%\opponentTwoJobPlace,\par
%        \opponentTwoJobPost
    \ifdefined\opponentThreeFio
        \par
        \vspace{0.01\paperheight}
        \textbf{\opponentThreeFio,}\par
        \opponentThreeRegalia,\par
        \opponentThreeJobPlace,\par
        \opponentThreeJobPost
    \fi
    }%
    \vspace{0.013\paperheight} \\
    \ifdefined\leadingOrganizationTitle
    \textbf{Ведущая организация:}    &
    \ifnumequal{\value{showopplead}}{0}{\vspace{6\onelineskip plus1fill}}{%
%        \leadingOrganizationTitle УДАЛИЛ ПОД ЕВСЮКОВА
		 \hspace{3cm}
    }%
    \fi
\end{tabularx}
\vspace{0.008\paperheight plus1fill}

\noindent Защита состоится \defenseDate~на~заседании диссертационного совета \defenseCouncilNumber~в \thesisInOrganization~по адресу: \defenseCouncilAddress.

\vspace{0.008\paperheight plus1fill}
\noindent Телефон для справок: \defenseCouncilPhone.

\vspace{0.008\paperheight plus1fill}

\noindent Ваш отзыв на автореферат в двух экземплярах, заверенных печатью, просим направлять по указанному адресу

%\noindent Отзывы на автореферат в двух экземплярах, заверенные печатью учреждения, просьба направлять по адресу: \defenseCouncilAddress, ученому секретарю диссертационного совета~\defenseCouncilNumber.

\vspace{0.008\paperheight plus1fill}
\noindent С диссертацией можно ознакомиться в библиотеке \thesisWhereOrganization.


\vspace{0.008\paperheight plus1fill}
\noindent{Автореферат разослан \synopsisDate.}



\vspace{0.008\paperheight plus1fill}
\noindent%
\begin{tabularx}{\textwidth}{@{}%
>{\raggedright\arraybackslash}b{18em}@{}
>{\centering\arraybackslash}X
r
@{}}
    Ученый секретарь\par
    диссертационного совета \defenseCouncilNumber\par
%    \defenseCouncilNumber,\par
    \defenseSecretaryRegaliaFull
    &
    \ifnumequal{\value{showsecrsign}}{0}{}{%
%        \includegraphics[width=2cm]{secretary-signature.png}%
    }%
    &
    \defenseSecretaryFioShort
\end{tabularx}

\vfill
\clearpage
           % Титульный лист
\include{Dissertation/contents}        % Оглавление
\ifnumequal{\value{contnumfig}}{1}{}{\counterwithout{figure}{chapter}}
\ifnumequal{\value{contnumtab}}{1}{}{\counterwithout{table}{chapter}}
\include{Dissertation/introduction}    % Введение
\ifnumequal{\value{contnumfig}}{1}{\counterwithout{figure}{chapter}
}{\counterwithin{figure}{chapter}}
\ifnumequal{\value{contnumtab}}{1}{\counterwithout{table}{chapter}
}{\counterwithin{table}{chapter}}
\include{Dissertation/part1}           % Глава 1
\include{Dissertation/part2}           % Глава 2
\include{Dissertation/part3}           % Глава 3
\chapter{Доработки для ученого совета МГТУ}\label{ch:ch4}

\section{Оформление разделов и подразделов}\label{sec:ch4/sect1}

Отделение точкой номеров разделов и подразделов
%Dissertation/setup.tex#L43:

%\setcounter{headingdelim}{2}


\subsection{Подпараграф \cyrdash{} один}\label{subsec:ch4/sect4/sub1}



%%%%%%%%%%%%%%%%%%%%%%%%%%%%%%%%%%%%%%%%%%%%%%%%%%%%%%%%%%%%%%%%%

LaTeX-шаблон для русской кандидатской диссертации и её автореферата.

%## Особенности
%* Кодировка: UTF-8.
%* Стандарт: ГОСТ Р 7.0.11-2011.
%* Поддерживаемые движки: pdfTeX, XeTeX, LuaTeX.
%* Поддерживаемые реализации библиографии: встроенная на движке BibTeX, BibLaTeX
%на движке Biber.

%[**Примеры компиляции шаблона**](https://github.com/AndreyAkinshin/Russian-Phd-LaTeX-Dissertation-Template/releases/latest).

%[**Установка программного обеспечения и сборка диссертации в файлы PDF**](Readme/Installation.md).

%[**Как писать диссертацию на GitHub?**](Readme/github.md)

Обсуждение

Общие вопросы лучше всего писать в gitter-канал:
%[![Join the chat at https://gitter.im/AndreyAkinshin/Russian-Phd-LaTeX-Dissertation-Template](https://badges.gitter.im/Join%20Chat.svg)](https://gitter.im/AndreyAkinshin/Russian-Phd-LaTeX-Dissertation-Template?utm_source=badge&utm_medium=badge&utm_campaign=pr-badge&utm_content=badge)

%Для отчётов об ошибках и для конкретных пожеланий/предложений лучше всего использовать раздел [Issues](https://github.com/AndreyAkinshin/Russian-Phd-LaTeX-Dissertation-Template/issues).

 Структура

%* [dissertation.tex](dissertation.tex): главный файл диссертации.
%* **[папка Dissertation](Dissertation/):** Структурированная система файлов с
шаблоном диссертации.
%  * **папка images:** Папка для размещения файлов изображений, относящихся только
%  к диссертации.
%  * [setup.tex](Dissertation/setup.tex): Файл упрощённой настройки оформления
%  диссертации.
%* [synopsis.tex](synopsis.tex): главный файл автореферата диссертации.
%* **[папка Synopsis](Synopsis/):** Структурированная система файлов с шаблоном
%автореферата.
%  * **папка images:** Папка для размещения файлов изображений, относящихся
%  только к автореферату диссертации.
%  * [setup.tex](Synopsis/setup.tex): Файл упрощённой настройки оформления
%  автореферата.
%* [presentation.tex](presentation.tex): главный файл презентации.
%* **[папка Presentation](Presentation/):** Структурированная система файлов с
%шаблоном презентации.
%  * **папка images:** Папка для размещения файлов изображений, относящихся
%  только к презентации.
%  * [setup.tex](Presentation/setup.tex): Файл упрощённой настройки оформления
%  презентации.
%* **[папка Documents](Documents/):** Полезные документы (ГОСТ-ы и постановления).
%* **[папка PSCyr](PSCyr/):** Пакет PSCyr + инструкции по установке.
%* **[папка BibTeX-Styles](BibTeX-Styles/):** Подборка русских стилевых пакетов
%BibTeX под UTF-8.
%* **[папка common](common/):** Общие файлы настроек и управления содержанием шаблонов.
%  * [characteristic.tex](common/characteristic.tex): Часть общей характеристики
%  работы, повторяющаяся в диссертации и автореферате.
%  * [concl.tex](common/concl.tex): Заключение. Является общим для автореферата
%  и диссертации (согласно [ГОСТ Р 7.0.11-2011](Documents/GOST%20R%207.0.11-2011.pdf),
%  пункты 5.3.3 и 9.2.3).
%  * [data.tex](common/data.tex): Общие данные (название работы, руководитель,
%  оппоненты, ключевые слова и т. п.).
%  * [packages.tex](common/packages.tex) и [styles.tex](common/styles.tex): Общие
%  пакеты и стили оформления автореферата и диссертации.
%  * [setup.tex](common/setup.tex): Общие настройки автореферата и диссертации.
%  В нём же настраивается выбор реализации библиографии.
%* **[папка biblio](biblio/):** Файлы с библиографией.
%  * [author.bib](biblio/author.bib): Публикации автора по теме диссертации.
%  * [registered.bib](biblio/registered.bib): Зарегистрированные патенты и программы для ЭВМ.
%  * [external.bib](biblio/external.bib): Работы которые ссылается автор.
%* **папка images:** Общие файлы изображений шаблонов.
%  * **папка cache:** Папка прекомпелированных рисунков.
%	* [placeholder.txt](images/cache/placeholder.txt): Файл, необходимый для прекомпиляции
%      рисунков в [overleaf](https://www.overleaf.com/).
%* **папка listings:** Общие файлы листингов.
%* **папка letters:** Файлы генерации конвертов для рассылки автореферата.

Дополнительные файлы:

%* [Makefile](Makefile), [compress.mk](compress.mk), [unix.mk](unix.mk),
%  [windows.mk](windows.mk), [examples.mk](examples.mk), [latexmkrc](latexmkrc): Файлы системы сборки шаблона.
%* [usercfg.mk](usercfg.mk): Пользовательские настройки системы сборки шаблона.
%* [indent.yaml](indent.yaml): Файл настройки форматирования исходного кода для
%  [latexindent](https://www.ctan.org/pkg/latexindent).
%* [.editorconfig](.editorconfig): Файл настройки текстовых редакторов, поддерживающих стандарт
%  [editorconfig](https://editorconfig.org/).
%* [Dockerfile](Dockerfile), [install-dockertex.sh](install-dockertex.sh): Файлы генерации
%  [Docker](https://www.docker.com/) образа для сборки шаблона.
%* [siunitx.cfg](siunitx.cfg): Определения величин SI для библиотеки
%  [siunitx](https://ctan.org/pkg/siunitx).
%* [synopsis_booklet.tex](synopsis_booklet.tex),
% [presentation_booklet.tex](presentation_booklet.tex): Файлы генерации печатных версий
% автореферата и презентации.
%* [presentation_handout.tex](presentation_handout.tex): Файл генерации раздаточных материалов из
%  презентации с добавлением подписей под слайдами.
%* [tikz.tex](tikz.tex): Файл изолированной сборки векторной графики [tikz](https://www.ctan.org/pkg/pgf).

 Дополнительная полезная информация

%* [Оформление библиографии](Readme/Bibliography.md)
%* [Как вносить правки в проект](CONTRIBUTING.md)
%* [Полезные ссылки](Readme/Links.md)
%* [Шаблон в галерее шаблонов ShareLaTeX](https://www.sharelatex.com/templates/thesis/russian-phd-latex-dissertation-template) (очень старая версия).

Благодарности

%* Большое спасибо Юлии Мартыновой за [оригинальный вариант шаблона](http://alessia-lano.livejournal.com/4267.html).
%* Большое спасибо [dustalov](https://github.com/dustalov),
%[Lenchik](https://github.com/Lenchik), [tonkonogov](https://github.com/tonkonogov)
%за значительный вклад и обсуждения.
%* Спасибо [storkvist](https://github.com/storkvist), [kshmirko](https://github.com/kshmirko),
%[ZoomRmc](https://github.com/ZoomRmc), [tonytonov](https://github.com/tonytonov),
%[Thibak](https://github.com/Thibak), [eximius8](https://github.com/eximius8),
%[Nizky](https://github.com/Nizky) за полезные правки и замечания.

Лицензия

%CC BY 4.0

Поэтому можно модифицировать и использовать шаблон любым образом, при
условии сохранения авторства на шаблон оформления диссертации в
формате LaTeX (в виде списка авторов в настоящем файле).  При этом не
накладывается никаких ограничений на текст диссертации, все права на
содержательную часть диссертации остаются за её автором.  В том числе,
если в тексте возникает раздел благодарностей (например, научному
руководителю за умелое руководство, коллегам за помощь в работе и
т.д.), то надо ли выносить авторам шаблона
*Russian-Phd-LaTeX-Dissertation-Template* благодарность за помощь в
оформлении диссертации или нет - решает сам диссертант. Использвание
шаблона не накладывает никаких ограничений на использование итоговых
файлов (например, PDF с готовой диссертацией или авторефератом),
т.е. никак не регулирует то, как они распространяются, копируются,
модифицируются и т.д.










\clearpage
           % Глава 4
\include{Dissertation/conclusion}      % Заключение
\include{Dissertation/acronyms}        % Список сокращений и условных обозначений
\include{Dissertation/dictionary}      % Словарь терминов
\include{Dissertation/references}      % Список литературы
\include{Dissertation/lists}           % Списки таблиц и изображений (иллюстративный материал)

\setcounter{totalchapter}{\value{chapter}} % Подсчёт количества глав

%%% Настройки для приложений
\appendix
% Оформление заголовков приложений ближе к ГОСТ:
\setlength{\midchapskip}{20pt}
\renewcommand*{\afterchapternum}{\par\nobreak\vskip \midchapskip}
\renewcommand\thechapter{\Asbuk{chapter}} % Чтобы приложения русскими буквами нумеровались

\include{Dissertation/appendix}        % Приложения

\setcounter{totalappendix}{\value{chapter}} % Подсчёт количества приложений

\end{document}
